\documentclass[12pt, a4paper]{article}
\usepackage[utf8]{inputenc}
\usepackage[T1]{fontenc}
\usepackage[margin=2.5cm]{geometry}
\usepackage{amsmath}
\usepackage{amssymb}
\usepackage{enumitem}
\usepackage{xcolor}
\usepackage{mdframed}
\usepackage{booktabs}
\usepackage{parskip}
\usepackage{titlesec}
\usepackage{hyperref}

\definecolor{truegreen}{RGB}{39, 80, 10}
\definecolor{truebg}{RGB}{234, 243, 222}
\definecolor{falsered}{RGB}{121, 31, 31}
\definecolor{falsebg}{RGB}{252, 235, 235}
\definecolor{notebg}{RGB}{230, 241, 251}
\definecolor{noteborder}{RGB}{24, 95, 165}
\definecolor{headergray}{RGB}{60, 60, 60}

\titleformat{\section}{\large\bfseries\color{headergray}}{}{0em}{}[\titlerule]
\titleformat{\subsection}{\normalsize\bfseries}{Q\thesubsection.}{0.5em}{}

\newcommand{\truelabel}{\textcolor{white}{\colorbox{truegreen}{\textbf{TRUE}}}}
\newcommand{\falselabel}{\textcolor{white}{\colorbox{falsered}{\textbf{FALSE}}}}

\newenvironment{statement}[1]{%
  \par\medskip
  \noindent #1\par\smallskip
  \noindent
}{\par}

\newenvironment{explanation}{%
  \noindent\small\color{headergray}\textit{Explanation: }%
}{\par}

\newenvironment{notebox}{%
  \begin{mdframed}[backgroundcolor=notebg, linecolor=noteborder, linewidth=1.5pt, leftline=true, rightline=false, topline=false, bottomline=false, innerleftmargin=10pt, innerrightmargin=10pt, innertopmargin=8pt, innerbottommargin=8pt]
}{\end{mdframed}}

\begin{document}

\begin{center}
  {\LARGE \textbf{Principal Component Analysis (PCA)}}\\[0.4em]
  {\large Exam Questions --- Answers \& Explanations}\\[0.2em]
  \textcolor{gray}{\small True/False with reasoning}
\end{center}

\vspace{1em}

\begin{notebox}
\textbf{Key Rules to Memorize}
\begin{itemize}[leftmargin=1.5em, itemsep=3pt]
  \item \textbf{Eigenvalues are always $\geq 0$.} Covariance matrices are positive semi-definite, so all eigenvalues are non-negative real numbers with no upper bound.
  \item \textbf{Sum of eigenvalues = total variance (trace), not generalised variance.} The generalised variance is the \textit{determinant} of the covariance matrix --- a completely different quantity.
  \item \textbf{Standardizing $\equiv$ using the correlation matrix.} Standardizing each variable to unit variance transforms the covariance matrix into the correlation matrix, so PCA on standardized data is identical to PCA on the correlation matrix.
  \item \textbf{Standardizing does not make eigenvalues equal.} It removes scale differences, but if variables are correlated, eigenvalues will still differ. Equal eigenvalues only arise when all variables are completely uncorrelated.
  \item \textbf{Total variance of standardized $p$-dimensional data $= p$.} Each standardized variable has variance 1, so the total is $p$, not 1.
  \item \textbf{PC scores are linear combinations of the original variables.} A score for PC$_i$ is $\mathbf{x}^\top \mathbf{w}_i$, a weighted sum of the original variables.
  \item \textbf{PCs are ordered by variance explained.} By construction, PC$_1$ captures the most variance, then PC$_2$, and so on: $\lambda_1 \geq \lambda_2 \geq \cdots \geq \lambda_p \geq 0$.
  \item \textbf{Eigenvectors of a covariance matrix are orthonormal.} A covariance matrix is symmetric, so its eigenvectors are orthogonal; after normalizing to unit length they form an orthonormal basis.
\end{itemize}
\end{notebox}

\vspace{1em}

%--------------------------------------------------------------
\section{Questions \& Answers}
%--------------------------------------------------------------

\subsection{PCA on a dataset with $p$ variables}

\textit{Which of the following statements are true?}

\bigskip

\begin{statement}{\truelabel\quad The first PC explains a proportion of total variance $\geq$ that of any of the $p$ original variables.}
\begin{explanation}
PC$_1$ is constructed to \emph{maximise} explained variance. In the standardized case each variable explains exactly $1/p$ of the total, whereas PC$_1$ always captures at least that much (and typically much more).
\end{explanation}
\end{statement}

\begin{statement}{\truelabel\quad The first PC explains a proportion of total variance $\geq$ that of any other principal component.}
\begin{explanation}
By definition, PCs are ranked in descending order of variance: $\lambda_1 \geq \lambda_2 \geq \cdots \geq \lambda_p$. So PC$_1$ always leads.
\end{explanation}
\end{statement}

\begin{statement}{\truelabel\quad If variables are standardized (mean 0, variance 1), they all contribute equally to total variance.}
\begin{explanation}
Standardizing sets each variable's variance to 1. The total variance equals $p$, and every variable contributes exactly $1/p$ of it --- equal contributions by construction.
\end{explanation}
\end{statement}

\begin{statement}{\falselabel\quad The sum of eigenvalues of the covariance matrix equals the generalised variance.}
\begin{explanation}
The sum of eigenvalues equals the \textbf{total variance} (the trace of the covariance matrix). The \textbf{generalised variance} is the \emph{determinant} of the covariance matrix --- a fundamentally different quantity.
\end{explanation}
\end{statement}

%--------------------------------------------------------------
\subsection{Typical range of eigenvalue values in PCA}

\textit{What is the typical range of values for eigenvalues in PCA?}

\bigskip

\begin{statement}{\falselabel\quad Non-negative real numbers, multiples of the variance.}
\begin{explanation}
Eigenvalues are non-negative real numbers, but there is no requirement that they be multiples of any specific variance value.
\end{explanation}
\end{statement}

\begin{statement}{\falselabel\quad Integers.}
\begin{explanation}
Eigenvalues are generally continuous real numbers. Nothing in the mathematics forces them to be whole numbers.
\end{explanation}
\end{statement}

\begin{statement}{\truelabel\quad Non-negative real numbers.}
\begin{explanation}
Covariance matrices are positive semi-definite, which guarantees $\lambda_i \geq 0$ for all $i$. Eigenvalues can take any non-negative real value.
\end{explanation}
\end{statement}

\begin{statement}{\falselabel\quad Real numbers between 0 and 1.}
\begin{explanation}
Eigenvalues have no upper bound of 1. For example, if a variable has very high variance, its associated eigenvalue could be 5, 100, or any large positive number.
\end{explanation}
\end{statement}

\begin{statement}{\falselabel\quad Real numbers (including negative).}
\begin{explanation}
Eigenvalues of a covariance matrix cannot be negative. The matrix is positive semi-definite by construction, ruling out negative eigenvalues.
\end{explanation}
\end{statement}

%--------------------------------------------------------------
\subsection{About PCA --- variance and eigenvalues}

\textit{Which of the following are true?}

\bigskip

\begin{statement}{\truelabel\quad Total variance is distributed among PCs proportionally to the eigenvalues of the covariance matrix.}
\begin{explanation}
The variance of PC$_i$ equals its eigenvalue $\lambda_i$. The proportion of variance it explains is $\lambda_i / \sum_j \lambda_j$, so eigenvalues directly determine the distribution of total variance.
\end{explanation}
\end{statement}

\begin{statement}{\falselabel\quad Total variance is distributed proportionally to the number of strictly positive loadings.}
\begin{explanation}
The distribution of variance depends entirely on eigenvalues, not on how many loadings happen to be positive. The sign and count of positive loadings carry no information about variance proportions.
\end{explanation}
\end{statement}

\begin{statement}{\falselabel\quad The eigenvalues of the covariance matrix can never all be equal.}
\begin{explanation}
If all $p$ variables are mutually uncorrelated and have the same variance $\sigma^2$, the covariance matrix is $\sigma^2 \mathbf{I}$ and all eigenvalues equal $\sigma^2$. Equal eigenvalues are uncommon but perfectly valid.
\end{explanation}
\end{statement}

\begin{statement}{\falselabel\quad If variables are standardized, the eigenvalues of their sample covariance matrix are all the same.}
\begin{explanation}
Standardizing removes scale differences but leaves the correlation structure intact. If variables are correlated, the eigenvalues will still be unequal. Equal eigenvalues require zero correlation among all variables.
\end{explanation}
\end{statement}

%--------------------------------------------------------------
\subsection{General properties of PCA}

\textit{Which statements are true regarding PCA?}

\bigskip

\begin{statement}{\truelabel\quad PCA provides uncorrelated linear combinations of the original variables.}
\begin{explanation}
Principal components are constructed to be mutually orthogonal (and hence uncorrelated). This is one of PCA's defining properties.
\end{explanation}
\end{statement}

\begin{statement}{\falselabel\quad Standardizing before PCA ensures that all eigenvalues are equal.}
\begin{explanation}
Standardizing changes the scale but not the correlation structure. Correlated variables produce unequal eigenvalues regardless of whether data were standardized first.
\end{explanation}
\end{statement}

\begin{statement}{\truelabel\quad The total variance explained by the first $k$ PCs equals the sum of the $k$ largest eigenvalues.}
\begin{explanation}
Since the variance of PC$_i$ is exactly $\lambda_i$, the cumulative variance of the top $k$ components is $\sum_{i=1}^{k} \lambda_i$.
\end{explanation}
\end{statement}

\begin{statement}{\truelabel\quad The eigenvectors of the sample covariance matrix form an orthonormal basis.}
\begin{explanation}
A covariance matrix is real and symmetric. The spectral theorem guarantees that such matrices have orthogonal eigenvectors. Normalizing them to unit length produces an orthonormal basis for $\mathbb{R}^p$.
\end{explanation}
\end{statement}

%--------------------------------------------------------------
\subsection{Scaling before PCA}

\textit{Which statements are correct regarding scaling before PCA?}

\bigskip

\begin{statement}{\truelabel\quad Without scaling, variables with large variance dominate the first PCs.}
\begin{explanation}
PCA maximises variance, so a variable expressed in millimetres (large numerical values) will overwhelm one expressed in kilometres. Scaling levels the playing field so all variables compete fairly.
\end{explanation}
\end{statement}

\begin{statement}{\truelabel\quad Centering without scaling is appropriate if variables are on the same scale.}
\begin{explanation}
When all variables share the same units and comparable magnitudes, centering (subtracting the mean) is sufficient. Full standardization would discard meaningful differences in spread.
\end{explanation}
\end{statement}

\begin{statement}{\truelabel\quad Standardizing is equivalent to applying PCA to the correlation matrix.}
\begin{explanation}
Standardizing each variable to unit variance converts the covariance matrix into the correlation matrix (since $\text{Cor}(X_i, X_j) = \text{Cov}(X_i/\sigma_i,\, X_j/\sigma_j)$). PCA on standardized data is therefore identical to PCA on the correlation matrix.
\end{explanation}
\end{statement}

\begin{statement}{\falselabel\quad PCA is unaffected by whether data are scaled or not.}
\begin{explanation}
Scaling fundamentally changes the analysis. Different scales produce a different covariance matrix, which yields different eigenvectors (principal components) and different eigenvalues (variance proportions).
\end{explanation}
\end{statement}

%--------------------------------------------------------------
\subsection{PCA on standardized 10-dimensional data}

\textit{Which of the following statements are true?}

\bigskip

\begin{statement}{\falselabel\quad The first PC corresponds to the direction of maximum correlation.}
\begin{explanation}
PC$_1$ is the direction of \textbf{maximum variance}, not maximum correlation. Correlation measures the linear relationship between two variables and is not equivalent to a direction in the principal component space.
\end{explanation}
\end{statement}

\begin{statement}{\falselabel\quad The loadings matrix is symmetrical.}
\begin{explanation}
The loadings matrix ($p \times p$) contains eigenvectors as its columns. In general it is \emph{not} symmetric. The covariance matrix is symmetric, but the matrix of its eigenvectors typically is not.
\end{explanation}
\end{statement}

\begin{statement}{\truelabel\quad The scores are linear combinations of the original variables.}
\begin{explanation}
The score for PC$_i$ is computed as $\mathbf{x}^\top \mathbf{w}_i$, which is a weighted sum (linear combination) of the original $p$ variables, with the loading vector $\mathbf{w}_i$ providing the weights.
\end{explanation}
\end{statement}

\begin{statement}{\falselabel\quad The total variance is equal to 1.}
\begin{explanation}
With 10 standardized variables, each has variance 1. The total variance is therefore $10 = p$, not 1. In general, total variance after standardization equals the number of variables $p$.
\end{explanation}
\end{statement}

\end{document}
